\subsection{}\etiquette[1.22]{VI.subsec1.22}
{\textit{\textbf{Exemples.}}} Reprenons l'exemple
\hyperref[VI.subsubsec1.6.2]{1.6.2} de la catégorie $(\Sch)$ avec une topologie
$T_i$, qui est un $𝒱$-site pour $𝒱$ convenable,
admettant la sous-catégorie des schémas affines comme catégorie
génératrice topologique formée d'objets quasi-compacts. On laisse au
lecteur le soin de vérifier qu'une flèche $f: X \to Y$ de ce site
est quasi-compacte (\resp quasi-séparé) si et seulement si c'est un
morphisme quasi-compact (\resp quasi-séparé) de schémas au sens
habituel (\cite{3} ou \cite{4}); de même un objet du site est
quasi-compact (\resp quasi-séparé) si et seulement si c'est un
schéma quasi-compact (\resp quasi-séparé) au sens habituel de
loc.cit. Il y a lieu de dire d'ailleurs, comme ici, \emph{morphisme
cohérent de schémas, schémas
  cohérent}, au lieu de « morphisme quasi-compact et quasi-séparé de
schémas », « schéma quasi-compact et quasi-séparé », comme on le
fait d'ailleurs dans \EGA I $2^{\text{ème}}$ édition.

\setcounter{subsection}{22}
\subsubsection{}\etiquette[1.22.1]{VI.subsubsec1.22.1}

Pour un schéma donné $X \in 𝒰$, on peut aussi regarder le topos
$\Top(X)$ associé à l'espace topologique sous-jacent, ou le topos
$\Top (X_{\et})$ (\resp $\Top(X_{\text{fppf}})$), associé au site des
$X$-schémas étales (\resp localement\pageoriginale de présentation
finie) éléments de $𝒰$, avec la topologie étale (\resp fppf). C'est un
$𝒰$-topos qui admet une famille génératrice formée d'objets cohérents,
correspondants aux objets affines du site de définition;

Ceci dit, l'objet final de ce topos est quasi-compact (\resp
quasi-séparé, \resp cohérent) si et seulement si le schéma $X$ est
quasi-compact (\resp quasi-séparé, \resp cohérent). De même, un
morphisme du site des espaces étalés sur $X$ (\resp du site
$X_{\text{et}}$ des schémas étales sur $X$, \resp du site
$X_{\text{fppf}}$ des schémas localement de présentation finie sur
$X$) est quasi-compact (\resp quasi-séparé, \resp cohérent) si et
seulement si c'est un morphisme de schémas qui est (au sens habituel
de \loccit) quasi-compact (\resp quasi-séparé, \resp cohérent).

\begin{enonce*}{Théorème 1.23}\etiquette[1.23]{VI.thm1.23}
Soient $E$ un $𝒰$-topos, $X$ un objet de $E$.
\begin{enumerate}
\renewcommand{\theenumi}{\roman{enumi}}
\renewcommand{\labelenumi}{(\theenumi)}
\item Pour que $X$ soit quasi-compact, il faut et il suffit que pour
  tout système inductif filtrant $(Y_i)_{i \in I}$ de $E$, avec $I \in
  𝒰$, l'application naturelle
\begin{equation*}
\varinjlim  \Hom(X, Y_i) \to \Hom(X, \varinjlim Y_i)\tag{1.23.1}
\end{equation*}\etiquette[1.23.1]{VI.eq1.23.1}
soit injective.

\item Supposons que $E$ admette une sous-catégorie génératrice $C$
  formée d'objets quasi-compacts. Pour que $X$ soit cohérent, il faut
  que pour toute donnée comme dans \emph{(i)}, l'application
  \emph{(1.23.1)} soit bijective.
\end{enumerate}
\end{enonce*}

\begin{enumerate}
\renewcommand{\theenumi}{\roman{enumi}}
\renewcommand{\labelenumi}{(\theenumi)}
\item \underline{Nécessité}.
Supposons\pageoriginale $X$ quasi-compact, il faut prouver que
  si $i\in I$ et $f_i$, $g_i : X \rightrightarrows Y_i$ sont tels que
  les composés $f$, $g$ de $f_i$, $g_i$ avec $Y_i \to Y = \varinjlim
  Y_j$ sont égaux, alors il existe $j \geq i$ tel que les composés
  $f_j$, $g_j$ avec $Y_i \to Y_j$ sont déjà égaux. Or comme dans $E$
  les limites inductives filtrantes commutent aux limites projectives
  finies, et en particulier aux noyaux, on a
$$
\Ker (f, g) = \mathop{\lim}_{j\geq l} \Ker (f_j, g_j){\quoi}.
$$
Par hypothèse $\Ker(f, g) = X$, donc $X$ est la limite de la famille
filtrante croissante de ses sous-objets $\Ker(f_j, g_j)$, donc comme
$X$ est quasi-compact, il est égal à un de ces sous-objets, donc il
existe $j \geq i$ tel que $f_j = g_{j}$.

\underline{Suffisance}. Pour prouver que $X$ est quasi-compact, il revient au
même (\hyperref[VI.subsubsec1.5.2]{1.5.2}) de prouver que pour toute famille filtrante croissante
$(X_i)_{i \in I}$ de sous-objets de $X$ dont la $\varinjlim$ est $X$,
il existe $i \in I$ tel que $X_i = X$. Or comme dans $E$ tout
monomorphisme est effectif (II 4.0), il s'ensuit que
$$
X_i = \Ker (f_i, g_i : X \rightrightarrows X \amalg_{X_{i}} X){\quoi}.
$$
Soit $Y_i = X \amalg_{X_i} X$, et considérons le système inductif formé
par les $Y_i$. Si $Y$ en est la limite inductive, on a $Y \simeq X
\amalg_{\varinjlim X_i} X = X \amalg_X X = X$.
Par suite, pour un $i \in I$ fixé, $f_i$, $g_i$ donnent le
même élément de $\Hom(X, Y)$, donc par hypothèse il existe $j \geq i$
tel que $f_j = g_j$ \ie $X_j = X$, \cqfd

\item Il reste\pageoriginale à prouver que sous les conditions indiquées, tout
  morphisme $X \to Y$ provient d'un morphisme $X \to Y_i$ pour un $i
  \in I$ convenable. Comme $C$ engendre $E$ et que les $Y_i$ couvrent
  $Y$, nous pouvons couvrir $X$ par des objets $X_{\alpha}$ de $C$, de
  telle façon que chacun des composés $X_{\alpha} \to X \to Y$ se
  factorise par un des $Y_i$. Comme $X$ est quasi-compact, on peut
  supposer la famille des $X_{\alpha}$ finie, donc que les morphismes
  $X_{\alpha} \to Y$ se factorisent par un $Y_i$ fixe en des
  $g_{\alpha i} : X_{\alpha} \to Y_i$. Comme $X$ est quasi-séparé, les
  $\fprod{X_{\alpha}}{X_{\beta}}{X}$ sont quasi-compacts. D'ailleurs
  pour tout $(\alpha, \beta)$, les composés de $\pr_1$ et $\pr_2$ sur
  $\fprod{X_{\alpha}}{X_{\beta}}{X}$ avec $X_{\alpha}$, $X_{\beta} \to
  Y_i$ sont
\[
\xymatrix{
\fprod{X_{\alpha}}{X_{\beta}}{X} \ar[rd] \ar[dd]& &\\
                            &X_{\alpha} \ar[ld] \ar[rrrr]&   & & & Y_i
\ar[lld]\\
X    \ar[rrr]                       &          & & Y  &
}
\]
tels que leurs composés avec $Y_i \to Y$ sont égaux. En
  vertu de(i), il existe donc un $j \geq i$ tel que les composés avec
  $Y_i \to Y_j$ soient déjà égaux. Comme l'ensemble des couples
  $(\alpha, \beta)$ est fini, on peut prendre $j$ indépendant de ce
  couple. Par suite les morphismes $g_{\alpha j}$ proviennent par
  composition d'un morphisme $g_j : X \to Y_j$, dont le composé avec
  $Y_j \to Y$ est d'ailleurs le morphisme $X \to Y$ donné, puisqu'il
  en est ainsi après composition avec les $Y_j \to X$. Cela achève la
  démonstration.
\end{enumerate}

\begin{enonce*}[remark]{Remarques 1.23.2}\etiquette[1.23.2]{VI.rem1.23.2}
L'hypothèse sur $E$ dans (ii) peut être remplacée par une hypothèse
supplémentaire \emph{sur} $X$, savoir que $X$ admet un système fondamental de
familles couvrantes par des objets quasi-compacts.
\end{enonce*}

\begin{enonce*}{Corollaire 1.24}\etiquette[1.24]{VI.cor1.24}
Soient\pageoriginale $E$ un topos, et soit $C$ une sous-catégorie
strictement pleine
de $E$ formée d'objets cohérents. Utilisant le fait que dans $E$ les
$𝒰$-limites inductives sont représentables, on trouve \emph{(I
  8.6.1)} un foncteur canonique.
\begin{equation*}
\Ind (C) \to E{\quoi}.\tag{1.24.1}
\end{equation*}\etiquette[1.24.1]{VI.eq1.24.1}
Ce foncteur est pleinement fidèle. Si $E$ admet une sous-catégorie
génératrice formée d'objets cohérents, et si $C$ est stable par
facteurs directs, alors les conditions suivantes sont équivalentes:
\begin{enumerate}
\renewcommand{\theenumi}{\roman{enumi}}
\renewcommand{\labelenumi}{(\theenumi)}
\item Le foncteur \emph{(1.24.1)} est essentiellement surjectif, \ie c'est
  une équivalence de catégories.

\item $C$ est une sous-catégorie génératrice de $E$, stable par
  limites inductives finies.

\item La catégorie $C$ est égale à la sous-catégorie pleine $E_{\coh}$
  de $E$ formée de tous les objets cohérents de $E$, et la condition
  nécessaire \emph{\hyperref[VI.thm1.23]{1.23} (ii)} de cohérence pour un objet de $E$ est
  aussi suffisante.
\end{enumerate}
\end{enonce*}

Soit $E_{PF}$ la sous-catégorie pleine de $E$ définie dans I 8.7.6. En
vertu de I 8.7.5 a), l'énoncé \hyperref[VI.thm1.23]{1.23} (ii) équivaut à la deuxième
inclusion de
$$
C \subset E_{\coh} \subset E_{PF}{\quoi},
$$
et l'inclusion composée $C \supset E_{PF}$ signifie que (1.24.1) est
pleinement fidèle. Comme\pageoriginale par hypothèse la sous-catégorie
$E_{\coh}$ de
$E$ est génératrice, il en est a fortiori de même de $E_{PF}$; d'autre
part dans $E$ les limites projectives finies sont représentables, de
sorte que nous pouvons appliquer I 8.7.7. La condition (iii) de \hyperref[VI.cor1.24]{1.24}
s'écrit aussi $C = E_{\coh} = E_{PF}$ et équivaut à la condition $C =
E_{PF}$, qui n'est autre que la condition (ii bis) de \loccit (compte
tenu que par hypothèse,  $E_{\coh}$ donc $E_{PF}$ est une
sous-catégorie génératrice de $E$). De même la condition (ii) de \hyperref[VI.cor1.24]{1.24}
n'est autre que la condition (iii bis) de \loccit. Donc l'équivalence
des conditions (i), (ii bis) et (iii bis) dans \loccit donne
l'équivalence des conditions (i), (ii) et (iii) de \hyperref[VI.cor1.24]{1.24}, \cqfd

La démonstration (ou plus précisément, l'invocation de I 8.7.7 (ii
bis)) fournit également la partie a) du

\begin{enonce*}{Corollaire 1.24.2}\etiquette[1.24.2]{VI.cor1.24.2}
Soit $E$ un topos admettant une sous-famille génératrice formée
d'objets cohérents, et soit $E_{PF}$ la sous-catégorie strictement
pleine des objets $X$ de $E$ tels que le foncteur covariant $\Hom(X,
-)$ qu'il représente commute aux limites inductives filtrantes.
Alors:
\begin{enumerate}
\renewcommand{\theenumi}{\alph{enumi}}
\renewcommand{\labelenumi}{\theenumi)}
\item Le foncteur naturel
$$
\Ind (E_{PF}) \to E
$$
est une équivalence de catégories.

\item Pour qu'un objet $X$ de $E$ appartienne à $E_{PF}$, il faut et
  il suffit qu'il soit isomorphe à un facteur direct ($=$ image d'un
  projecteur (I, 10.6)) du conoyau\pageoriginale d'une double flèche $X_1
  \rightrightarrows X_0$, avec $X_1$ et $X_0$ cohérents.
\item La sous-catégorie $E_{PF}$ de $E$ est stable par limites
  inductives finies, et est équivalente à une petite catégorie.
\end{enumerate}
\end{enonce*}

Il reste à prouver b) et c). La première assertion de c) est triviale
grâce à I 2.8, et implique le « il suffit » de b) grâce à \hyperref[VI.thm1.23]{1.23}
(ii). Pour le « il faut », on utilise \hyperref[VI.thm1.23]{1.23} (i) qui montre que $X$
est quasi-compact, d'où l'existence d'un épimorphisme $X_{0} \to X$, avec
$X_0$ cohérent, compte tenu du fait que $E_{\coh}$ engendre $E$ et est
stable par sommes finies. Soit $R = \fprod{X_0}{X_0}{X}$, de sorte que
$X$ s'identifie au conoyau de $R \rightrightarrows X_0$. En vertu de
a), on a $R = \varinjlim R_i$, limite inductive filtrante des coker
$(R_i \rightrightarrows X_0) = X_i$. D'après l'hypothèse $X \in \Ob
E_{PF}$, pour $i$ assez grand le morphisme $u_i : X'_i \to X$ admet
une section $v : X \to X'_i$. Soit $w = vu_i : X'_i \to
X'_i$, de sorte que $w^2 = w$ et $X \simeq \Coker (w, \id_{X'_i}
: X'_i \rightrightarrows X'_i)$. Comme $R_i$ est à nouveau
quasi-compact, on le couvre par $X_{1} \to R_i$ avec $X_i$
cohérent, d'où $X_i \simeq \Coker (X_1 \rightrightarrows X_0)$, ce qui
prouve b).

Enfin, comme $E_{PF}\subset E_{\text{qu.cpct}}$, la dernière assertion de c)
résulte de \hyperref[VI.subsubsec1.5.3]{1.5.3}.

\begin{enonce*}{Corollaire 1.25}\etiquette[1.25]{VI.cor1.25}
Soit $E$ un topos tel que la sous-catégorie pleine $C$ de $E$ formée
des objets cohérents soit génératrice. Considérons le foncteur
canonique
$$
\varphi : \Ind(c) \to E{\quoi}.
$$
Les conditions\pageoriginale suivantes sont équivalentes:
\begin{itemize}
\item[(i)] Le foncteur $\varphi$ est essentiellement surjectif, \ie tout
  objet de $E$ est limite inductive filtrante d'objets cohérents.

\item[(i bis)] Le foncteur $\varphi$ est une équivalence de catégories.

\item[(ii)] Toute limite inductive finie dans $E$ d'objets cohérents est
  un objet cohérent.

\item[(ii bis)] Le conoyau dans $E$ d'une double flèche d'objets
  cohérents est un objet cohérent.

\item[(iii)] La réciproque de \emph{\hyperref[VI.thm1.23]{1.23} (ii)} est valable.
\end{itemize}
\end{enonce*}

C'est un cas particulier de \hyperref[VI.cor1.24]{1.24}, compte tenu que $C$ est stable par
sommes finies (\hyperref[VI.cor1.4]{1.4}), ce qui implique l'équivalence de (ii) et de (ii
bis), et compte tenu du

\begin{enonce*}{Lemme 1.25.1}\etiquette[1.25.1]{VI.lem1.25.1}
Soit $E$ un topos. Tout facteur direct $X'$ d'un objet cohérent $X$ de
$E$ est cohérent.
\end{enonce*}

En effet, $X'$ est quasi-compact comme quotient de $X$ (\hyperref[VI.prop1.3]{1.3}), et
quasi-séparé comme sous-objet de $X$ (\hyperref[VI.cor1.15]{1.15}).

\begin{enonce*}{Corollaire 1.26}\etiquette[1.26]{VI.cor1.26}
Soit $E$ un topos satisfaisant aux conditions équivalentes
\emph{\hyperref[VI.cor1.25]{1.25}.} Alors, il en est de même de tout topos induit.
\end{enonce*}

Cela résulte aussitôt du critère (ii), compte tenu de \hyperref[VI.subsubsec1.20.2]{1.20.2} et du
fait que l'existence d'une sous-catégorie génératrice de $E$ formée
d'objets cohérents implique manifestement la même propriété pour un
topos induit.

\begin{enonce*}[remark]{Remarque 1.27}\etiquette[1.27]{VI.rem1.27}
Parmi\pageoriginale les exemples importants de topos satisfaisant les
conditions de \hyperref[VI.cor1.25]{1.25}, signalons les topos noethériens
(\hyperref[VI.subsec2.14]{2.14}), et le topos zariskien ou étale
(\hyperrefext[VII][VII.sec1]{1}) d'un schéma $X$ (\cf IX, note p.42). Mais
même lorsque
le topos $E$ est \emph{cohérent}(\hyperref[VI.def2.3]{2.3}) (\ie la sous-catégorie pleine
$C$ des objet cohérents est génératrice et stable par limites
projectives finies), il n'est par toujours vrai que $C$ soit
\emph{parfait} (\hyperref[VI.def2.9.1]{2.9.1}) \ie satisfasse aux conditions équivalentes de
\hyperref[VI.cor1.25]{1.25}; \cf \hyperref[VI.exer1.28]{1.28} ci-dessous.
\end{enonce*}

\begin{enonce*}[remark]{Exercice 1.28}\etiquette[1.28]{VI.exer1.28}
Soient $E$ un topos, $p_1$, $p_2 : X_1 \rightrightarrows X_0$ une double
flèche dans $E$. Définir une suite croissante de sous-objets $R_n$ $(n
\geq 0)$ de $X_0\times X_0$ par la condition $R_0 = \Sup(\text{Im}(p_1,
p_2), \text{Im}(p_2, p_1))$, et $R_n = \pr_{13} (\pr_{12}^{-1} (R_{n-1}),
\pr_{23}^{-1}(R_{n-1}))$ pour tout $n \geq 1$, où $\pr_{ij} (1 \leq i
< j \leq 3)$ désignent les projections qu'on devine du produit triple
de $X_0$ vers le produit double.
\begin{enumerate}
\renewcommand{\theenumi}{\alph{enumi}}
\renewcommand{\labelenumi}{\theenumi)}
\item Montrer qu'on obtient ainsi une suite croissante de sous-objets
  de $X_0\times X_0$, que $R = \varinjlim R_n$ est un graphe
  d'équivalence (I 10.4) dans $X_0$, et que le morphisme canonique
$$
X = \Coker (p_1, p_2) \to X_0 /R
$$
est un isomorphisme.

\item En conclure que si $X_0$ est quasi-compact et $X = \Coker(p_1,
  p_2)$ est quasi-séparé (donc cohérent), alors la suite des $R_n$ est
  stationnaire. Inversement,\pageoriginale si cette suite est
  stationnaire, et si on suppose $X_0$
cohérent et $X_1$ quasi-compact alors $X$ est cohérent. (Pour ce
dernier point, écrire $R_n \simeq \fprod{(R_{n-1}, \pr_2)}{(R_{n-1},
  \pr_1)}{X_0}$ et en conclure que $R_n$ est quasi-compact sur $X_0$
pour tout $n$, puis utiliser 1.19.)

\item Construire un exemple de deux applications d'ensembles $p_1$, $p_2
  : X_1 \rightrightarrows X_0$, telles que la suite des $(R_n)$ ne
  soit pas stationnaire.

\item Soit $C$ une petite catégorie où les limites inductives finies
  et les limites projectives finies sont représentables, et où tout
  morphisme se factorise en un épimorphisme effectif suivi d'un
  monomorphisme. Répéter pour une double flèche $(p_1, p_2)$ de $C$ la
  construction des $R_n$. Munissant $C$ de la topologie canonique,
  montrer que le foncteur canonique $\epsilon : C \to C^{\sim}$ est exact, et
  commute par suite à la formation des $R_n$.

\item Prendre pour $C$ une petite sous-catégorie pleine de $(\Ens)$,
  stable à isomorphisme près par limites projectives finies et limites
  inductives finies, et contenant les ensembles $X_0$, $X_1$ de c). En
  conclure que le topos $C^{\sim}$ (qui est un topos
  \emph{cohérent} (\hyperref[VI.def2.3]{2.3}), \ie engendré par la sous-catégorie de ses
  objets cohérents, laquelle est stable par limites projectives finies)
  ne satisfait pas aux conditions équivalentes de \hyperref[VI.cor1.25]{1.25}.

\item Soient $k$ un corps, $C$ le site fppf des schémas localement de
  présentation finie sur $k$ (\SGA 3 IV 6.3). Montrer qu'il existe un
  schéma affine réduit $X$ sur $k$ qui admit un automorphisme $f$ qui
  n'est pas d'ordre fini (prendre par exemple l'espace affine $E^2$
  muni de l'automorphisme $f(x) = x$, $f(y) = x + y$).\pageoriginale
  Montrer que si
   $p_1$, $p_2 : X_1 \rightrightarrows X_0$ est une double flèche dans
  $C$, telle que la double flèche correspondant dans $(\Ens) p_1
  (\bar{k}):p_2(\bar{k}): X_1(\bar{k}) \rightrightarrows X_0
  (\bar{k})$ donne une suite $(R_n)$ non stationnaire, alors il en
  est de même de la double flèche de $C^{\sim}$ définie par $p_1$,
  $p_2$, donc, si $X_0$ est de type fini sur $k$, alors le conoyau
  dans $C^{\sim}$ de $(p_1, p_2)$ n'est pas cohérent. En conclure que
  le conoyau dans $C^{\sim}$ du couple $(f, \id_X)$ n'est pas
  cohérent, donc que le topos $C^{\sim}$ ne satisfait par aux
  conditions équivalentes de \hyperref[VI.cor1.25]{1.25}. En conclure plus généralement que
  si $S$ est un schéma non vide, $C$ le site fppf des schémas
  localement de présentation finie sur $S$, alors le topos $C^{\sim}$
  ne satisfait pas aux conditions de \hyperref[VI.cor1.25]{1.25}. (On fera attention que,
  contrairement à l'apparence, $C^{\sim}$ n'est donc noethérien (\hyperref[VI.def2.11]{2.11})
  que si $S$ est vide, comme il résulte de ce qui précède et de \hyperref[VI.subsec2.14]{2.14}.
\end{enumerate}
\end{enonce*}

\begin{enonce*}{Définition 1.30}\etiquette[1.30]{VI.def1.30}
Soit $E$ un topos. Un objet $X$ de $E$ est dit un \emph{objet
prénoethérien}\footnote{Par la notion plus forte d'objet
  \emph{noethérien}, \cf \hyperref[VI.def2.11]{2.11} ci-dessous.} du topos $E$ s'il satisfait
aux deux conditions équivalentes suivantes:
\begin{enumerate}
\renewcommand{\theenumi}{\roman{enumi}}
\renewcommand{\labelenumi}{(\theenumi)}

\item Tout sous-objet de $X$ est quasi-compact.

\item Toute suite croissante de sous-objets de $X$ est stationnaire.
\end{enumerate}
\end{enonce*}

\setcounter{subsection}{30}
\setcounter{subsubsection}{0}
\subsubsection{}\etiquette[1.30.1]{VI.subsubsec1.30.1}
L'équivalence des conditions (i) et (ii) est claire. On notera que ces
conditions ne dépendent encore que du topos induit $E_{/X}$, et même
seulement de l'ensemble ordonné des ouverts de $E_{/X}$, isomorphe à
l'ensemble ordonné des sous-objets de $X$. Elles sont stables par
passage à un sous-objet de $X$.


\subsection{}\etiquette[1.31]{VI.subsec1.31}
On prouve\pageoriginale comme pour \hyperref[VI.prop1.3]{1.3} et
\hyperref[VI.cor1.4]{1.4} les faits suivantes : si $(X_i \to X)$ est une
famille couvrante finie, avec les $X_i$ prénoethériens, alors $X$
est prénoethérien; en particulier un quotient d'un objet
prénoethérien de $E$ est prénoethérien. Si $(X_i)_{i \in I}$ est une
famille d'objets de $E$, alors leur somme $X$ est un objet
prénoethérien de $E$ si et seulement si tous les $X_i$ sauf un
nombre fini sont isomorphes à $∅_{E}$, et tous les $X_i$
sont prénoethériens.

\subsection{}\etiquette[1.32]{VI.subsec1.32}
Évidemment un objet prénoethérien d'un topos $E$ est quasi-compact;
lorsque $E$ admet une famille génératrice formée d'objets
prénoethériens, la réciproque est vraie (comme il résulte aussitôt
de \hyperref[VI.subsec1.31]{1.31}): les objets prénoethériens de $E$ sont alors
ses objets quasi-compacts.

Supposons que $E$ admette une famille génératrice $(X_i)_{i\in I}$
formée d'objets quasi-compacts, alors les trois conditions suivantes
sont équivalentes:
\begin{enumerate}
\renewcommand{\theenumi}{\roman{enumi}}
\renewcommand{\labelenumi}{(\theenumi)}
\item Les $X_i$ sont prénoethériens.

\item Tout objet quasi-compact de $E$ est prénoethérien.

\item Tout monomorphisme dans $E$ est quasi-compact.
\end{enumerate}
\noindent
(On vient de voir l'équivalence de (i) et (ii), (ii) $\Rightarrow$
(iii) est trivial à partir des définitions, et (iii) $\Rightarrow$
(ii) sur la définition \hyperref[VI.def1.30]{1.30} (i).)

On notera\pageoriginale que (iii) a comme conséquence que \emph{tout
  morphisme de
  $E$ est quasi-séparé}; donc (\hyperref[VI.prop1.14]{1.14} (ii)) tout objet $X$ de $E$
  au-dessus d'un objet quasi-séparé $X$ de $E$ est lui-même
  quasi-séparé.

\begin{enonce*}[remark]{Exemple 1.33}\etiquette[1.33]{VI.exem1.33}
(\emph{Topos classifiant d'un groupe}). Soient $G$ un groupe $\in
𝒰$, $B_G$ son topos classifiant (IV 2.4) \ie la catégorie
des ensemble $\in 𝒰$ où $G$ opère à gauche. Les objets
connexes-non vides de $B_G$ (IV 4.3.5) sont les ensembles à
opérateurs qui sont non vides et sur lesquels $G$ opère
transitivement, et tout objet $E$ de $B_G$ est somme de ses
composantes connexes, qui sont les sous-objets non vides minimaux de
$E$ (savoir les orbites de $G$ dans $E$). Par suite $E$ est un objet
quasi-compact de $B_G$ si et seulement si l'ensemble $\pi_0(E)$ des
orbites de $G$ dans $E$ est fini, et alors $E$ est même un objet
prénoethérien de $B_G$. Ces objets (et même le seul objet $G_s$, qui
est connexe non-vide) forment une sous-catégorie génératrice de
$B_G$. Pour que l'objet final $e$ de $B_G$ soit quasi-séparé, \ie
pour que le produit de deux objets quasi-compacts de $B_G$ soit
quasi-compact, il faut et il suffit que le produit $G_s\times G_s$
soit quasi-compact, ce qui équivaut manifestement au fait que $G$
est fini. On en conclut, plus généralement, qu'un objet $E$ de $B_G$
est quasi-séparé si et seulement si les stabilisateurs de ses points
sont des sous-groupes \emph{finis} de $G$; \ie s'il est isomorphe à
une somme d'espaces homogènes $G/H$, avec $H$ \emph{fini}.) En
effet, en vertu de \hyperref[VI.cor1.15]{1.15} on peut supposer dans ce critère
$E$ connexe, mais si $x \in E$ a comme groupe de stabilité $G_x$, on
sait (IV 5.8)\pageoriginale que $B_G / E$ est équivalent au topos
$B_{G_x}$, donc son objet final est quasi-séparé si et seulement si
$G_x$ est fini. On conclut de ce critère que tout objet de $B_G$ qui
se trouve au-dessus d'un objet quasi-séparé est quasi-séparé, a
fortiori, le produit de deux objets quasi-séparés de $B_G$ est
quasi-séparé. On conclut aussi que les objets cohérents de $B_G$
sont les objets isomorphes à des sommes finies d'espaces homogènes
$G/H$, avec $H$ \emph{fini}. Comme $G_s$ est cohérent, on voit que
les objets cohérents de $B_G$ forment déjà une sous-catégorie
génératrice de $B_G$. De plus, cette sous-catégorie est stable par
produits fibrés (comme il résulte de fait que tout objet au-dessus
d'un objet quasi-séparé est quasi-séparé). (Dans la terminologie
\hyperref[VI.def2.11]{2.11} plus bas $B_G$ est un topos \emph{localement
noethérien}, mais il n'est quasi-séparé que si $G$ est fini, et
alors $B_G$ est même noethérien.)

On voit tout de suite qu'un morphisme $f : E'\to E$ est
quasi-compact si et seulement si l'application induit $\pi_0(f) :
\pi_0 (E')\to \pi_{0}(E)$ pour les ensembles d'orbites est propre
\ie à fibres finies; on retrouve qu'un monomorphisme est
quasi-compact (\hyperref[VI.subsec1.32]{1.32}), donc tout morphisme est
quasi-séparé. Par suite les morphismes cohérents dans $B_G$ sont
simplement les morphismes quasi-compacts, qu'on vient de
caractériser.

Soit $E$ un objet du topos $B_G$. On vérifie facilement que $E$ est
une limite inductive filtrante d'objets cohérents de $B_G$ si et
seulement si les stabilisateurs des points de $E$ sont des groupes
ind-finis (\ie limites\pageoriginale inductives filtrantes de leurs
sous-groupes finis). En particulier, si $G$ n'est pas ind-fini (par
exemple si $G =  \ZZ$), l'objet final du topos $B_G$ n'est pas
limite inductive d'objets cohérents; plus précisément, pour que le topos $B_G$
satisfasse aux conditions équivalentes de \hyperref[VI.cor1.25]{1.25}, il faut et il suffit
que son objet final soit limite inductive filtrante d'objets
cohérents, ou encore que le groupe $G$ soit ind-fini.
\end{enonce*}


\section{Conditions de finitude pour un topos}\etiquette[2]{VI.sec2}

\begin{enonce*}{Proposition 2.1}\etiquette[2.1]{VI.prop2.1}
Soit $E$ un $𝒰$-topos. Les conditions suivantes sont équivalentes
\begin{enumerate}
\renewcommand{\theenumi}{\roman{enumi}}
\renewcommand{\labelenumi}{(\theenumi)}
\item Il existe une sous-catégorie pleine génératrice $C$ de $E$,
  formée d'objets quasi-compacts, et stable par produits fibrés.

\item Il existe un $𝒰$-site $C$, tel que tout objet de $C$ soit
  quasi-compact, que dans $C$ les produits fibrés soient
  représentables, et que $E$ soit équivalent à la catégorie des
  faisceaux $C^{\sim}$.
\end{enumerate}

Supposons que ces conditions sont satisfaites. Alors on peut même
choisir dans \emph{(i) (\resp (ii))} la catégorie $C$
$𝒰$-petite. D'autre part, pour tout $X \in \ob C$, l'objet
$X$ \emph{(\resp $\epsilon (X)$)} du topos $E$ est \emph{cohérent}.
\end{enonce*}

L'implication (i) $\Rightarrow$ (ii) résulte de IV 1.2.1., et le
fait que dans (i) (\resp (ii)) on peut prendre $C$ $𝒰$-petite se voit
aussitôt par l'argument de IV 1.2.3 appliqué à une petite
sous-catégorie pleine génératrice au sens topologique (II 3.0.1) de
$C$. Il reste à prouver (ii) $\Rightarrow$ (i) et la dernière
assertion de \hyperref[VI.prop2.1]{2.1}, ce qui résulte aussitôt du

\begin{enonce*}{Corollaire 2.1.1}\etiquette[2.1.1]{VI.cor2.1.1}
Soit\pageoriginale $C$ un $𝒰$-site où les produits fibrés soient
représentables et dont tout objet soit quasi-compact, et soient $E : C^{\sim}$,
$\varepsilon : C \to E$ le foncteur canonique. Alors pour tout $X \in
\ob C$, $\varepsilon(X)$ est un objet cohérent de $E$. La
sous-catégorie pleine de $E$ formée des objets cohérents de $E$ qui se
trouvent au-dessus de quelque $\varepsilon(X)$ est stable par produits
fibrés (et est évidemment génératrice dans $E$).
\end{enonce*}

On sait déjà que $\epsilon(X)$ est quasi-compact (\hyperref[VI.prop1.2]{1.2}),
prouvons qu'il est quasi-séparé. Soient donc $F$, $G$ deux objets
quasi-compacts de $E$ et $F \to \varepsilon(X)$, $G \to \epsilon(X)$
deux morphismes, prouvons que le produit fibré
$\fprod{F}{G}{\epsilon(X)}$ est quasi-compact. Recouvrant $F$ et $G$
par un nombre fini d'objets de la forme $\epsilon(Y_i)$ \resp
$\epsilon(Z_j)$, et procédant comme dans \hyperref[VI.prop1.2]{1.2} pour nous
ramener au cas où les composés $\epsilon(Y_i) \to \epsilon(X)$,
$\epsilon(Z_j) \to \epsilon(X)$ proviennent de morphismes $Y_i \to
X$ \resp $Z_j \to Y$, on est ramené grâce à \hyperref[VI.prop1.3]{1.3} au cas
où les morphismes $F \to \epsilon(X)$ et $G \to \epsilon(X)$
sont les transformés par $\epsilon$ de morphismes $Y \to X$, $Z \to
X$. Mais comme les produits fibrés sont représentables dans $C$ et
que $\epsilon$ y commute, on est réduit à prouver que
$\epsilon(\fprod{Y}{Z}{X})$ est quasi-compact, ce qui a été déjà
noté plus haut pour tout objet de la forme $\epsilon(U)$.

Soient maintenant $G \to F$, $H \to F$ des morphismes d'objets
cohérents de $E$, tels qu'il existe un morphisme $F \to
\epsilon(X)$, pour quelque $X \in \ob C$, prouvons que
$\fprod{G}{H}{F}$ est cohérent. Comme $G$ et $H$ sont quasi-compacts
et $F$ quasi-séparé, il résulte déjà des définitions que le produit
fibré est quasi-compact.\pageoriginale Reste à voir qu'il est
quasi-séparé, ou ce qui revient au même (\hyperref[VI.cor1.15]{1.15}) que $\fprod{G}{H}{\epsilon(X)}$ est
quasi-séparé. Cela résulte du fait plus général:

\begin{enonce*}{Corollaire 2.1.2}\etiquette[2.1.2]{VI.cor2.1.2}
Sous les conditions de \hyperref[VI.cor2.1.1]{2.1.1}, pour tout objet quasi-séparé $G$ de $E$,
tout morphisme $G \to \epsilon(X)$ est quasi-séparé. Si $H \to
\epsilon(X)$ est un deuxième morphisme, avec $H$ quasi-séparé,
alors $\fprod{G}{H}{\epsilon(X)}$ est quasi-séparé.
\end{enonce*}

La première assertion implique la deuxième, car elle implique par
changement de base que $\fprod{G}{H}{\epsilon(X)}$ est quasi-séparé,
donc, puisque $H$ est quasi-séparé, $\fprod{G}{H}{\epsilon(X)}$
l'est aussi (\hyperref[VI.prop1.14]{1.14} (ii)). Pour prouver que $G \to \epsilon(X)$ est quasi-séparé,
nous allons utiliser le

\begin{enonce*}{Lemme 2.1.3}\etiquette[2.1.3]{VI.lem2.1.3}
Soient $E$ un topos, $f : X \to Y$ un morphisme dans $E$. Supposons
que $E$ admette une famille génératrice d'objets quasi-compacts, et
qu'il existe une famille couvrante $(g_i : X_i \to X)_{i \in I}$ de
$X$ telle que les $X_i$ soient quasi-compacts, et que pour tout couple
d'indices $i$, $j \in I$, $\fprod{X_i}{X_j}{Y}$ soit
quasi-séparé. Alors, si $X$ est un objet quasi-séparé de $E$, $f$ est
un morphisme quasi-séparé.
\end{enonce*}

En effet, la famille $(\fprod{g_i}{g_j}{Y} : \fprod{X_i}{X_j}{Y} \to
\fprod{X}{X}{Y})_{(i, j) \in I \times I}$ est couvrante, donc pour
voir que le morphisme $\Delta_{X/Y} : X \to \fprod{X}{X}{Y}$ est
quasi-compact, il suffit de voir que pour tout $(i, j) \in I \times
I$, le morphisme qui s'en\pageoriginale déduit par changement de base
$\fprod{g_i}{g_j}{Y}$ l'est (\hyperref[VI.prop1.10]{1.10} (ii)). Or ce dernier n'est autre que
le morphisme canonique $\fprod{X_i}{X_j}{X} \to \fprod{X_i}{X_j}{Y}$;
comme par hypothèse $X$ est quasi-séparé et les $X_i$ sont
quasi-compacts, $\fprod{X_{i}}{X_{j}}{X}$ est quasi-compact, donc
comme $\fprod{X_{i}}{X_{j}}{Y}$ est quasi-séparé par hypothèse, le
morphisme $\fprod{X_{i}}{X_{j}}{X}\to \fprod{X_{i}}{X_{j}}{Y}$ est
bien quasi-compact,\cqfd

Nous pouvons maintenant terminer la démonstration de \hyperref[VI.cor2.1.2]{2.1.2}, en
prouvant que tout morphisme $f : G \to \epsilon(X)$, avec $G$ avec
$G$ quasi-séparé, est quasi-séparé. Pour ceci, couvrons $G$ par des
objets $\epsilon(X_i)$, tels que les morphismes composés
$\epsilon(X_i) \to \epsilon(X)$ proviennent de morphismes $X_i
\to X$. En vertu de \hyperref[VI.lem2.1.3]{2.1.3}, comme les $\epsilon(X_i)$ sont
quasi-compacts, il suffit de vérifier que les produits fibrés
$\fprod{\epsilon(X_i)}{\epsilon(X_j)}{\epsilon(X)} =
\epsilon(\fprod{X_i}{X_j}{X})$ sont quasi-séparés. Or on a déjà vu
que tout objet de $E$ de la forme $\epsilon(U)$ est cohérent. Cela
prouve \hyperref[VI.cor2.1.2]{2.1.2} et achève la démonstration de \hyperref[VI.prop2.1]{2.1}.

\begin{enonce*}{Proposition 2.2}\etiquette[2.2]{VI.prop2.2}
Soit $E$ un topos contenant une sous-catégorie pleine génératrice $C$
formée d'objets cohérents. Les conditions suivantes sur $E$ sont
équivalentes:
\begin{itemize}
\item[(i)] Tout objet quasi-séparé de $E$ est quasi-séparé sur l'objet
  final.

\item[(i bis)] Pour tout morphisme $f : X \to Y$ dans $E$, $X$
  quasi-séparé implique $f$ quasi-sépare.

\item[(i ter)] Tout objet $X$ de $C$ est quasi-séparé sur l'objet final de $E$,
  \ie le morphisme diagonal $X \to X \times X$ est quasi-compact.

\item[(ii)] Le produit\pageoriginale de deux objets quasi-séparés de $E$ est
  quasi-séparé.

\item[(ii bis)] Le produit dans $E$ de deux objets de $C$ est quasi-séparé.
\end{itemize}

Ces conditions impliquent que la sous-catégorie pleine $E_{\coh}$ de
$E$ formée des objets cohérents de $E$ est stable par produits fibrés
(et à fortiori, que $E$ satisfait aux conditions équivalentes de \hyperref[VI.prop2.1]{2.1}).
\end{enonce*}

Évidemment (i bis) $\Rightarrow$ (i), et l'implication inverse
résulte de \hyperref[VI.prop1.8]{1.8} (iii). On a (i) $\Rightarrow$ (ii) par
un argument déjà fait : si $X$ et $Y$ sont quasi-séparés, comme par
l'hypothèse (i) $X$ est quasi-séparé sur $e$, $X \times Y$ est
quasi-séparé sur $Y$ par changement de base, donc $X$ est
quasi-séparé en vertu de \hyperref[VI.prop1.14]{1.14} (ii). Le même argument
montre que (i ter) $\Rightarrow$ (ii bis), et d'autre part (i)
$\Rightarrow$ (i ter) et (ii) $\Rightarrow$ (ii bis) sont triviaux,
de sorte qu'il reste à établir (ii bis) $\Rightarrow$ (i). Mais si
$X$ est un objet quasi-séparé de $E$, recouvrant $X$ par des objet
$X_i$ de $C$ (ce qui est possible, $C$ étant génératrice), comme les
$X_i$ sont quasi-compacts par hypothèse sur $C$, on peut appliquer
\hyperref[VI.lem2.1.3]{2.1.3} au morphisme $X \to e$, de sorte que pour prouver
que ce dernier est quasi-séparé, on est ramené précisément à voir
que les $X_i \times X_j$ sont quasi-séparés, ce qui est l'hypothèse
(ii bis).

\begin{enonce*}[remark]{Remarque 2.2.1}\etiquette[2.2.1]{VI.rem2.2.1}
La condition (i ter) équivaut aussi à la condition (i quater). Pour
toute double flèche $g_1$, $g_2 : Y \rightrightarrows X$  dans $C$, le
noyau dans $E$ est quasi-compact (ou encore (\hyperref[VI.cor1.15]{1.15}) cohérent).
\end{enonce*}

Pour le\pageoriginale voir, il suffit d'appliquer le critère \hyperref[VI.prop1.10]{1.10} (i)
au morphisme diagonal $X \to X \times X$ envisagé dans (i ter).

\begin{enonce*}{Définition 2.3}\etiquette[2.3]{VI.def2.3}
Un topos $E$ satisfaisant les conditions équivalentes de \emph{\hyperref[VI.prop2.2]{2.2} (\resp de
\hyperref[VI.prop2.1]{2.1})} est appelé un \emph{topos algébrique (\resp} un \emph{topos localement
algébrique}, ou \emph{localement cohérent)}. Un objet $X$ d'un topos $E$ est
appelé \emph{objet algébrique du topos} $E$ si le topos induite $E_{/X}$ est
algébrique. Un topos $E$ est dit \emph{quasi-séparé (\resp cohérent)} s'il
est algébrique et si son objet final est quasi-séparé (\resp
cohérent).
\end{enonce*}

\setcounter{subsection}{3}
\subsubsection{}\etiquette[2.3.1]{VI.subsubsec2.3.1}
C'est dans les topos algébriques que les notions de finitude
développées au ${\rm n^{\circ}}$~1 ont des propriétés pleinement
satisfaisantes, analogues aux propriétés des notions de même nom dans
la catégorie des schémas. Tous les topos localement algébriques (=
localement cohérents) rencontrés en pratique sont en fait algébriques,
donc l'intérêt pratique de cette notion semble pour l'instant assez
réduit. On peut cependant construire des topos localement algébriques
et non algébriques (\hyperref[VI.exer2.17]{2.17} f).

\subsection{}\etiquette[2.4]{VI.subsec2.4}
Voici quelques remarques générales concernant les notions introduites
dans \hyperref[VI.def2.3]{2.3}, qui convaincront le lecteur (du moins on
l'espère) que la terminologie introduite est cohérente.

\subsubsection{}\etiquette[2.4.1]{VI.subsubsec2.4.1}
Sous les conditions de \hyperref[VI.prop2.1]{2.1} (i) \resp (ii), pour tout $X \in C$, $X$
(\resp $\epsilon(X)$) est un objet (cohérent) algébrique de $E$
(\hyperref[VI.cor2.1.2]{2.1.2} et 1.19.2),
en d'autres\pageoriginale termes le topos induit
$E_{/X}$ (\resp $E_{\epsilon(X)} \approxeq (C_{/X})^{\sim}$) est
algébrique, donc
\emph{cohérent}. Il revient au même de dire que \emph{lorsque
$C$ admet un objet
final, alors $E$ lui-même est algébrique, donc cohérent.}

\subsubsection{}\etiquette[2.4.2]{VI.subsubsec2.4.2}
Si $E$ est un topos algébrique, alors tout topos induit $E_{/X}$ est
également algébrique (\hyperref[VI.prop2.2]{2.2} (i bis)); pour qu'un topos $E$ soit
localement algébrique = localement cohérent, il faut et il suffit
qu'il existe une famille $(X_i)_{i \in I}$ d'objets de $X$, couvrant
l'objet final, telle que les $X_i$ soient algébriques (\resp
cohérents) \ie telle que les topos induits $E_{/X_{i}}$ soient algébriques
(\resp cohérents). -La suffisance résulte aussitôt des définitions, la
nécessité des observations \hyperref[VI.subsubsec2.4.1]{2.4.1}. Bien entendu, un topos algébrique
est localement algébrique, et si $E$ est un topos localement
algébrique, tout topos induit $E_{/X}$ est localement algébrique.

\subsubsection{}\etiquette[2.4.3]{VI.subsubsec2.4.3}
Si $X$ est un objet algébrique d'un topos $E$, alors tout objet $X'$
de $E$ qui se trouve au-dessus de $X$ est encore algébrique (\cf
\hyperref[VI.subsubsec2.4.2]{2.4.2}). En particulier, tout objet d'un topos algébrique est
algébrique, et inversement bien sûr.

\subsubsection{}\etiquette[2.4.4]{VI.subsubsec2.4.4}
La sous-catégorie pleine de $E$ formée des objets
\emph{quasi-séparés} qui sont \emph{algébriques} est stable par
produits fibrés et par sous-objets (\hyperref[VI.cor1.15]{1.15}), donc aussi
par noyaux. Si $E$ est algébrique, cette catégorie (qui n'est alors
autre que la catégorie des objets quasi-séparés de $E$
sans\pageoriginale plus) est également stable par produit de deux
objets. Si $E$ est quasi-séparé (et alors seulement) cette catégorie
est stable par limites projectives finis quelconques, ou encore
contient l'objet final de $E$.

La sous-catégorie pleine $C$ de $E$ formée des objets \emph{cohérents}
qui sont \emph{algébriques} est stable par produits fibrés. Si $E$ est
localement algébrique \ie localement cohérent, dire que $E$ est
algébrique revient à dire que $C$ est également stable par noyaux de
doubles flèches (\hyperref[VI.rem2.2.1]{2.2.1}); et dire que $E$ est quasi-séparé revient à
dire que de plus $C$ est stable par produit de deux objets (\cf
\hyperref[VI.cor1.17]{1.17}). Enfin dire que $E$ est cohérent revient à dire que $C$ est
stables par limites projectives finies quelconques, ou encore qu'il
contient l'objet final de $E$.

\subsubsection{}\etiquette[2.4.5]{VI.subsubsec2.4.5}
Soit $E$ un topos. Pour que $E$ soit localement cohérent (\resp
algébrique, \resp quasi-séparé, \resp cohérent) il faut et il suffit
qu'il admette une sous-catégorie pleine génératrice $C$ formée
d'objets quasi-compacts (qui seront automatiquement cohérents
(\hyperref[VI.prop2.1]{2.1})), et qui soit \emph{stable} par produits fibrés
(\resp par produits fibrés et par noyaux, \resp par produits fibrés
et par produits de deux objets, \resp par limites projectives finies
quelconques); ou encore que $E$ soit équivalent à un topos
$C^{\sim}$, où $C$ est un $𝒰$-site dont tout objet est
quasi-compact, et où les produits fibrés et les produits de deux
objets (\resp toutes les limites projectives finies) sont
représentables. (Pour la nécessité, on prend \pageoriginale pour $C$
la catégorie des objets de $E$ qui sont cohérents et
\emph{algébriques}, et on applique \hyperref[VI.subsec2.4]{2.4}; pour la
suffisance, on applique \hyperref[VI.prop2.2]{2.2} (i ter)). Dans cet énoncé,
on peut évidemment supposer encore $C$ $𝒰$-petit.

\subsubsection{}\etiquette[2.4.6]{VI.subsubsec2.4.6}
Soit $E$ un topos, Si $E$ est algébrique, un objet $X$ de $E$ est
quasi-séparé (\resp cohérent) si et seulement si le topos induit
$E_{/X}$ l'est. Lorsque $E$ n'est plus supposé algébrique, on peut
dire seulement que $E_{/X}$ est quasi-séparé (\resp cohérent) si et
seulement si $X$ est quasi-séparé (\resp cohérent) et
\emph{algébrique}. Ce n'est pas grave, vu qu'on n'aura sans doute pas
à utiliser ces notions en dehors du cas $E$ algébrique.

\subsubsection{}\etiquette[2.4.7]{VI.subsubsec2.4.7}
Pour que la notion d'objet cohérent (\resp quasi-séparé d'un
topos $E$ ait des propriétés satisfaisantes, il faut se borner aux
objets qui sont de plus algébriques, condition automatiquement
satisfaite si $E$ est algébrique. Comme nous n'aurons sans doute
jamais à travailler avec des $E$ localement algébriques qui ne soient
algébriques, la question de réviser la terminologie introduite dans
\hyperref[VI.def1.13]{1.13}, en la réservant éventuellement aux seuls objets algébriques, ne
se pose donc pas en termes bien aigus!

D'autre part, nous sommes abstenus de définir la notion de topos
quasi-compact. Dans l'esprit du présent numéro, il s'imposerait
d'appeler ainsi les topos \emph{algébriques} sont l'objet final est
quasi-compact.

Nous \pageoriginale avons hésité à introduire un tel usage, car si
$X$ est un espace topologique, il ne serait pas vrai que $X$ est
quasi-compact si et seulement si le topos associé $\Top(X)$ l'est.
On notera à ce propos que $\Top(X)$ est localement cohérent si et
seulement si $X$ admet une base d'ouverts formée d'ouverts
quasi-compacts $U$ tels que pour $U_j$, $U_k \subset U_i$, $U_j \cap
U_k$ soit encore quasi-compact (et alors $\Top(X)$ est même
algébrique). Cette conditions est vérifiée pour les espaces
topologiques sous-jacents aux schémas, mais rarement pour les
espaces rencontrés par les analystes, fussent-ils (les espaces)
compacts. Explicitons encore, pour la commodité des références:

\begin{enonce*}{Proposition 2.5}\etiquette[2.5]{VI.prop2.5}
Soit $E$ un topos algébrique, $f : X \to Y$ un morphisme dans $E$. Si
$Y$ est un objet quasi-séparé (\resp cohérent) de $E$, alors, pour que
$X$ soit un objet quasi-séparé (\resp cohérent) de $E$, il faut et il
suffit que le morphisme $f$ soit quasi-séparé (\resp cohérent).
\end{enonce*}

La suffisance a déjà été vue dans \hyperref[VI.prop1.14]{1.14}(ii), et est vraie sans
hypothèse sur $E$. La nécessité dans la cas non respé est vraie sans
supposer même $Y$ quasi-séparé, et n'est  autre que la définition \hyperref[VI.def2.3]{2.3}
via \hyperref[VI.prop2.2]{2.2} i bis). Il reste à prouver que si $Y$ est cohérent et $X$
cohérent, alors $f$ est cohérent. Comme on sait déjà que $f$ est
quasi-séparé, il suffit de voir que $f$ est quasi-compact, ce qui
résulte du fait que $X$ est quasi-compact et $Y$ quasi-séparé (\hyperref[VI.def1.13]{1.13}).

\begin{enonce*}{Corollaire 2.6}\etiquette[2.6]{VI.cor2.6}
Soient\pageoriginale $E$ un topos algébrique, $f: X \to Y$ un
morphisme dans $E$, $(Y_i \to Y)_{i \in I}$ une famille couvrante, avec les $Y_i$
quasi-séparés (\resp  cohérents). Pour que $f$ soit quasi-sépare
(\resp quasi-compact, \resp cohérent) il faut et il suffit que pour
tout $i \in I$, $X_i = \fprod{X}{Y_i}{Y}$ soit un objet quasi-séparé
(\resp quasi-compact, \resp cohérent) de $E$.
\end{enonce*}

Il suffit de conjuguer \hyperref[VI.prop1.10]{1.10} (ii) et \hyperref[VI.prop2.5]{2.5}
dans le cas « quasi séparé » ou « cohérent »; le cas « quasi-compact » a déjà été vu dans \hyperref[VI.cor1.16]{1.16}. On notera que
la nécessité de la condition énoncée dans \hyperref[VI.cor2.6]{2.6} est
également valable sans hypothèse sur le topos $E$.

On trouve comme cas particulier de \hyperref[VI.cor2.6]{2.6}:

\begin{enonce*}{Corollaire 2.7}\etiquette[2.7]{VI.cor2.7}
Soient $E$ un topos cohérent, $X$ un objet de $E$. Pour que $X$ soit
un objet quasi-compact (\resp quasi-séparé, \resp cohérent) de $E$, il
faut et il suffit qu'il soit quasi-compact (\resp quasi-séparé, \resp
cohérent) au-dessus de l'objet final (\hyperref[VI.subsubsec1.9.2]{1.9.2}).
\end{enonce*}

En particulier, $X$ est constructible (\hyperref[VI.subsubsec1.9.2]{1.9.2}) si et seulement s'il
est cohérent.

\begin{enonce*}{Corollaire 2.8}\etiquette[2.8]{VI.cor2.8}
Soient $E$ un topos localement cohérent, $f:X \to Y$ un morphisme dans
$E$. Pour que $f$ soit quasi-séparé, il faut et il suffit que pour
tout objet quasi-séparé $Y'$ de $E_{/Y}$, $f^{\ast} (Y') =
\fprod{X}{Y'}{Y}$ soit un objet quasi-séparé de $E_{/X}$ (ou de $E$,
cela revient au même (\hyperref[VI.subsubsec1.20.2]{1.20.2})).
\end{enonce*}

Le \pageoriginale « il faut » est valable sans condition sur $E$,
car si $Y'$ est quasi-séparé, comme $f':X'= \fprod{X}{Y'}{Y} \to Y'$ est quasi-séparé par
changement de base, $X'$ est quasi-séparé (\hyperref[VI.prop1.14]{1.14} (ii)). Pour la
réciproque, notons qu'il existe une famille couvrante $Y_i \to Y$ de
$Y$ par des objets quasi-séparés et algébriques. En vertu de \hyperref[VI.prop1.10]{1.10}
(ii), pour vérifier que $f$ est quasi-séparé, il suffit de vérifier
que les $f_i : X_i = \fprod{X}{X_i}{Y} \to Y_i$ le sont. Or
l'hypothèse sur $f^{\ast}$ implique que les $X_i$ sont quasi-séparés,
donc les $f_i$ le sont puisque $E_{/Y_i}$ est un topos algébrique.

\begin{enonce*}[remark]{Remarque 2.8.1}\etiquette[2.8.1]{VI.rem2.8.1}
Les énoncés \hyperref[VI.prop2.5]{2.5} et \hyperref[VI.cor2.6]{2.6} restent valables si on y remplace l'hypothèse
« $E$ algébrique » par l'hypothèse plus générale « $Y$ est
algébrique », comme on voit trivialement en appliquant l'énoncé
primitif au topos induit $E_{/Y}$, qui est algébrique. De même, dans
\hyperref[VI.cor2.8]{2.8} il suffit de supposer que $E_{/Y}$ (au lieu de $E$) soit
localement cohérent.
\end{enonce*}

\begin{enonce*}{Proposition 2.9}\etiquette[2.9]{VI.prop2.9}
Soient $E$ un topos, $C$ une sous-catégorie de $E$. Les conditions
suivantes sont équivalentes:
\begin{enumerate}
\renewcommand{\theenumi}{\roman{enumi}}
\renewcommand{\labelenumi}{(\theenumi)}
\item $E$ est localement cohérent (\resp cohérent), satisfait aux
  conditions équivalentes de \hyperref[VI.cor1.25]{1.25}, et $C$ est la sous-catégorie
  pleine $E_{\coh}$ de $E$ formée des objets cohérents.

\item $C$ est une sous-catégorie strictement pleine génératrice de
  $E$, stable par limites inductives finies et par produits fibrés
  (\resp et par limites projectives finies), et est formée d'objets
  quasi-compacts.
\end{enumerate}
\end{enonce*}

L'implication\pageoriginale (i) $\Rightarrow$ (ii) résulte trivialement des
définitions et de la forme \hyperref[VI.cor1.25]{1.25} (ii) des conditions envisagées dans
\hyperref[VI.cor1.25]{1.25}; prouvons l'implication inverse. Le fait que $C$ soit une
sous-catégorie pleine génératrice, formée d'objet quasi-compacts, et
stable par produits fibrés (\resp par limites projectives finies)
implique, par définition, que $E$ est localement cohérent (\resp
cohérent). Le fait que de plus $C$ soit strictement pleine et stable
par limites inductives finies implique alors que $C = E_{\coh}$ (en
vertu de \hyperref[VI.cor1.24]{1.24} (ii) $\Rightarrow$ (iii)), et que la condition \hyperref[VI.cor1.25]{1.25}
(ii) est satisfaite, \cqfd

\begin{enonce*}{Définition 2.9.1}\etiquette[2.9.1]{VI.def2.9.1}
Un $𝒰$-topos $E$ est dit parfait s'il est cohérent (\hyperref[VI.def2.3]{2.3}) et s'il
satisfait aux conditions équivalentes de \hyperref[VI.cor1.25]{1.25}. \ie si la
sous-catégorie $E_{\coh}$ de $E$ formée des objets cohérents de $E$ est
stable par limites inductives finies. On dit que $E$ est localement
parfait s'il existe une famille $(X_i)_{i \in I}$ d'objets de $E$
couvrant l'objet final, telle que pour tout $i \in I$, le topos induit
$E_{/X_i}$ soit parfait.
\end{enonce*}

\setcounter{subsection}{9}
\setcounter{subsubsection}{1}
\subsubsection{}\etiquette[2.9.2]{VI.subsubsec2.9.2}

Il résulte aussitôt de cette définition qu'un topos parfait (\resp
localement parfait) est cohérent (\resp localement cohérent). Comme
exemples de topos parfaits (\resp localement parfaits), signalons
les topos noethériens (\resp localement noethériens) introduits dans
\hyperref[VI.def2.11]{2.11} ci-dessous (\cf \hyperref[VI.subsec2.14]{2.14}), le topos
zariskien et le topos étale d'un schéma cohérent (\resp d'un schéma
quelconque) (IX, note page 42). Comme contre-exemple,\pageoriginale
signalons le topos fppf d'un schéma $S$, qui est localement cohérent
(et même cohérent si $S$ est un schéma cohérent, \ie quasi-compact
et quasi-séparé) (\hyperrefext[VII][VII.lem5.6]{5.6}), mais qui n'est pas localement
parfait si $S \neq ∅$ (\hyperref[VI.exer1.28]{1.28} f)).

Dans un topos localement parfait, la notion d'objet constructible
est particulièrement stable (\hyperref[VI.prop2.9.3]{2.9.3} ci-dessous), ce qui
est une raison pourquoi la notion semble intéressante. Une autre
raison est dans le fait que comme celle de topos cohérent ou
localement cohérent, la notion de topos parfait ou localement
parfait est stable par rapport à la formation du sous-topos fermé
complémentaire d'un ouvert (\hyperref[VI.prop4.6]{4.6}), et qu'elle se
comporte de façon particulièrement simple pour l'opération de
recollement de topos (\hyperref[VI.cor4.11]{4.11}, \hyperref[VI.cor4.14]{4.14}).

\begin{enonce*}{Proposition 2.9.3}\etiquette[2.9.3]{VI.prop2.9.3}
Soit $E$ un topos localement parfait (\hyperref[VI.def2.9.1]{2.9.1}). Alors la
sous-catégorie
pleine $E_{\cons}$ de $E$ formée des objets constructibles de $E$
(\hyperref[VI.subsubsec1.9.3]{1.9.3}) est stable par limites inductives finies (et
aussi par limites
projectives finies bien sûr, en vertu de (\hyperref[VI.subsubsec1.9.3]{1.9.3})).
\end{enonce*}

Comme la propriété pour un objet de $E$ d'être constructible est
locale sur $E$ (\hyperref[VI.prop1.10]{1.10} (ii)), on est ramené au cas
où $E$ est un topos
parfait. Mais alors $E_{\cons} = E_{\coh}$ (\hyperref[VI.cor2.7]{2.7}), et la conclusion
résulte de la définition \hyperref[VI.def2.9.1]{2.9.1}.

\begin{enonce*}{Corollaire 2.9.4}\etiquette[2.9.4]{VI.cor2.9.4}
Soit\pageoriginale $E$ un topos localement parfait. Alors $E$ est parfait si et
seulement si $E$ est cohérent. Pour tout objet $X$ de $E$, le topos
induit $E_{/X}$ est localement parfait.
\end{enonce*}

\begin{enonce*}[remark]{Problème 2.9.5}\etiquette[2.9.5]{VI.prob2.9.5}
Il est concevable que les topos parfaits soient assez particuliers
pour se prêter à une \emph{théorie de structure} aussi explicite
(suivant un modèle proposé par M.\textsc{Artin} à l'époque du
séminaire oral). Appelons \emph{topos fini} un topos équivalent à un
topos de la forme $\widehat{C}$, où $C$  est une catégorie finie
(\cf exercice \hyperref[VI.exer3.11]{3.11}), et \emph{topos profini} un topos
qui est une limite projective filtrante de topos finis (au sans de
6. plus bas, qui s'applique grâce à \hyperref[VI.exer3.12]{3.12} c)). Comme un
topos fini est noethérien (\hyperref[VI.exer2.17]{2.17} g)) donc parfait, et
qu'une limite projective filtrante de topos parfaits est parfait
(6.), on voit qu'un topos profini est parfait. La question qui se
pose serait de savoir si réciproquement tout topos parfait est
profini. Cela équivaut à la question si toute sous-catégorie finie
de $E_{\coh}$ est contenue dans une sous-catégorie pleine $F_{0}$ de
$E_{\coh}$ qui est équivalente à une catégorie $F_{\coh}$, où
$F$ est un topos fini (\cf \hyperref[VI.exer3.11]{3.11}). (Si la réponse était
négative, il y aurait lieu de trouver des conditions intrinsèques
supplémentaires maniables pour un topos parfait qui assurent qu'il
est profini.) Il resterait enfin à étudier la structure des topos
profinis en termes d'une notion convenable de « catégorie profinie
karoubienne », inspirée de IV 7.6 h). Cette\pageoriginale étude
n'a pas été faite encore même dans le cas particulier
où $E$ est
le topos zariskien, ou étale, d'un brave schéma cohérent $X$, et
devrait donner alors des invariants plus fins que l'espace compact
$X_{\cons}$ (\EGA IV 1.) associé à $X$.

Dans le même ordre d'idées, signalons la questions suivante: Soit
$E$ un topos parfait, $E'$ le topos induit sur un ouvert de $E$, et
$E''$ le sous-topos fermé correspondant (IV 9.9). Si $E'$ et $E''$
sont des topos finis, en est-il de même de $E$?
\end{enonce*}

\begin{enonce*}{Proposition 2.10}\etiquette[2.10]{VI.prop2.10}
Soit $E$ un topos. Les conditions suivantes sont équivalentes:
\begin{itemize}
\item[(i)] $E$ admet une sous-catégorie pleine génératrice
  stable par
  produits fibrés, formée d'objets prénoethériens de $E$.

\item[(ii)] $E$ admet une sous-catégorie pleine
  génératrice formée d'objets
  prénoethériens quasi-séparés (\ie prénoethériens cohérents).

\item[(iii)] $E$ est localement cohérent (\hyperref[VI.def2.3]{2.3}) et tout objet quasi-compact de
  $E$ est prénoethérien.

\item[(iii bis)] $E$ est localement cohérent et admet une famille
  génératrice formée d'objets prénoethériens de $E$.
\end{itemize}
\end{enonce*}

L'équivalence de (iii) et (iii bis) résulte de \hyperref[VI.subsec1.32]{1.32},
et il est trivial que (i) $\Rightarrow$ (iii bis) et (iii)
$\Rightarrow$ (i), donc (i) équivaut à (iii) et (iii bis). Enfin (i)
$\Rightarrow$ (ii) résulte de la dernière assertion de
\hyperref[VI.prop2.1]{2.1},\pageoriginale et il reste à prouver (ii)
$\Rightarrow$ (i). Cette implication résulte du fait que la
sous-catégorie pleine de $E$ formée des objets noethériens
quasi-séparés, si elle est génératrice, est stable par produits
fibrés, car un tel produit fibré est quasi-compact en vertu des
définitions, donc noethérien par la première assertion de
\hyperref[VI.subsec1.32]{1.32}; et il est quasi-séparé par la dernière
assertion de \hyperref[VI.subsec1.32]{1.32}.

\begin{enonce*}{Définition 2.11}\etiquette[2.11]{VI.def2.11}
Un topos $E$ est dit topos localement noethérien s'il satisfait aux
conditions équivalentes de \hyperref[VI.prop2.10]{2.10}, noethérien si de
plus son objet final
est cohérent (ou ce qui revient au même (\hyperref[VI.subsec1.32]{1.32}))
prénoethérien et
quasi-séparé). Un objet $X$ d'un topos $E$ est dit objet
noethérien de
$E$ si le topos induit $E_{/X}$ est un topos noethérien.
\end{enonce*}

\setcounter{subsection}{11}
\subsection{}\etiquette[2.12]{VI.subsec2.12}
Si $E$ est un topos localement noethérien, alors il en est de même de
tout topos induit $E_{/X}$; inversement, si on peut trouver une
famille $(X_i)$ couvrant l'objet final telle que les $E_{X_i}$ soient
localement noethériens, il en est de même de $E$. Si $C$ est une
sous-catégorie pleine génératrice de $E$ formée d'objets
prénoethériens et stable par produits fibrés (\hyperref[VI.def2.11]{2.11} (i)), alors pour
tout $X \in C$, $E_{/X}$ est un topos noethérien, \ie les objets de
$C$ sont même \emph{noethériens}, (car si $E$ est localement
noethérien et $X \in \Ob E$, alors $E_{/X}$ est noethérien si et
seulement si $X$ est cohérent). Il s'ensuit qu'un topos $E$ est
localement noethérien si et seulement si on peut recouvrir l'objet
final de $E$ par des objets noethériens $X_i$.

\subsection{}\etiquette[2.13]{VI.subsec2.13}
Dans un\pageoriginale topos localement noethérien $E$, tout monomorphisme est
quasi-compact, et tout morphisme est quasi-séparé (\hyperref[VI.subsec1.32]{1.32}). A fortiori,
$E$ est un topos \emph{algébrique} (\hyperref[VI.def2.3]{2.3}) et non seulement localement
cohérent \ie localement algébrique. Donc un topos $E$ est noethérien
si et seulement si il est localement noethérien et \emph{cohérent}
(\hyperref[VI.def2.3]{2.3}). En d'autres termes, si $E$ est un topos localement noethérien,
alors les objets noethériens de $E$ ne sont autres que les objets
\emph{cohérents} de $E$.

\subsection{}\etiquette[2.14]{VI.subsec2.14}
Soit $E$ un topos localement noethérien, Si $E$ est quasi-séparé,
\ie si son objet final est quasi-séparé, alors \emph{tout} objet de
$E$ est quasi-séparé (\hyperref[VI.subsec1.32]{1.32}), donc les objets
prénoethériens de $E$ sont identiques aux objets noethériens (\ie
cohérents) de $E$. Comme un objet quotient d'un objet prénoethérien
est prénoethérien, il s'ensuit que $E$ satisfait aux conditions
équivalentes de \hyperref[VI.cor1.25]{1.25} (puisqu'il satisfait la condition
\hyperref[VI.cor1.25]{1.25} (ii bis)). En particulier, si $C$ est la
sous-catégorie pleine de $E$ formée des objets noethériens
(cohérents) de $E$, alors le foncteur naturel
$$
\Ind(C) \to E
$$
est une équivalence de catégories. On conclut de ceci qu'\emph{un topos
noethérien est parfait} (\hyperref[VI.def2.9.1]{2.9.1}), donc qu'un topos localement
noethérien est localement parfait.

Si on suppose seulement $E$ localement noethérien, il sera encore
vrai que tout objet de $E$ est limite inductive de ses sous-objets
prénoethériens (car dans un topos admettant une sous-catégorie
génératrice formée\pageoriginale d'objets quasi-compacts, tout objet
est limite inductive filtrante de ses sous-objets quasi-compacts).
Mais comme un objet prénoethérien de $E$ n'est plus nécessairement
cohérent, cet énoncé n'a alors qu'une utilité très limitée, et on
sait (\hyperref[VI.exem1.33]{1.33}) que les conditions de \hyperref[VI.cor1.25]{1.25} ne
sont plus nécessairement vérifiées.

\begin{enonce*}[remark]{Exemple 2.15}\etiquette[2.15]{VI.exem2.15}
(\emph{Espaces topologiques}.) Soit $X$ un espace topologique. Alors
les conditions suivantes sont équivalentes : (i) l'espace
topologique $X$ est noethérien (\ie tout ouvert de $X$ est
quasi-compact), \ie toute suite croissante d'ouverts de $X$ est
stationnaire); (ii) le topos $\Top(X)$ est noethérien; (iii) l'objet
final $X$ de $\Top(X)$ est noethérien; (iv) l'objet final $X$ de
$\Top(X)$ est prénoethérien.

En effet, on a trivialement (ii) $\Leftrightarrow$ (iii) $\Rightarrow$
(iv) $\Leftrightarrow$ (i), d'autre part (i) implique que la
sous-catégorie génératrice ouv(X) de $\Top(X)$, qui est stable par
limites projectives finies, est formée d'objets prénoethériens, d'où (ii).

On voit de même que l'espace topologique $X$ est localement
noethérien (\ie est réunion d'ouverts noethériens) si et seulement
si le topos $\Top(X)$ est localement noethérien.
\end{enonce*}

\setcounter{subsection}{15}
\setcounter{subsubsection}{0}
\subsubsection{}\etiquette[2.15.1]{VI.subsubsec2.15.1}
Ainsi, nos définitions \hyperref[VI.def1.30]{1.30}, \hyperref[VI.def2.11]{2.11} sont compatibles avec la
terminologie reçue pour les espaces topologique. D'autre part, nous
avons tenu dans le cas général à donner au terme « objet
noethérien » un sens\pageoriginale plus fort que celui de la notion
plus naïve de \hyperref[VI.def1.30]{1.30}, pour que l'ensemble des propriétés qui s'attachent à cette
notion (plutôt que la seule structure grammaticale de la définition)
soit bien en accord avec l'intuition qui s'attache aux espaces
topologiques et aux schémas \emph{noethériens}.

Le lecteur trouvera d'autres exemples dans les exercices suivants.

\begin{enonce*}[remark]{Exercice 2.16}\etiquette[2.16]{VI.exer2.16}
(\emph{Espaces à opérateurs et topos classifiants}).

Soit $X$ un espace topologique sur lequel opère un groupe discret
$G$, d'où (IV 2.3) un topos $E = \Top(X, G)$. On interprète les
objets $X'$ de ce topos comme des espaces étalés sur $X$ à groupe
d'opérateurs $G$, et on note que le topos induit $E_{/X}$, est
canoniquement équivalent à $\Top(X', G)$.
\begin{enumerate}
\renewcommand{\theenumi}{\alph{enumi}}
\renewcommand{\labelenumi}{\theenumi)}
\item Pour que l'objet final du topos $E$ soit quasi-compact, il faut
  et il suffit que l'espace topologique quotient $X/G$ soit
  quasi-compact (Utiliser IV 8.4.1).

\item Soit $P$ l'objet de $E$ définit par l'espace étalé trivial $X
  \times G$ avec opération diagonale de $G$. Montrer que le topos
  induit $E_{/P}$ est équivalent au topos $\Top(X)$ (comparer IV
  5.8.3). En conclure que $E$ est localement cohérent (\resp
  localement noethérien) si et seulement si $\Top(X)$ est localement
  cohérent (\cf \hyperref[VI.subsubsec2.4.7]{2.4.7}) (\resp localement noethérien (\cf
  \hyperref[VI.exem2.15]{2.15})). Montrer que si $E$ est localement cohérent, \ie localement
  algébrique, il est même algébrique, en utilisant la famille
  génératrice formée\pageoriginale des objets $T_U$, où $U$ est un
  ouvert cohérent   de $X$, $T_U = G \times U$ avec opération $g \cdot (u, g') =
  (u, gg')$ de $G$, et la morphisme structural $p_U : T_U \to X$ défini
  par $P_U(u, g) = g\cdot u$.

\item Supposons $E$ algébrique, \ie $\Top(X)$ algébrique
  (\hyperref[VI.subsubsec2.4.7]{2.4.7}). Montrer que le morphisme structural $P \to e$ est
  quasi-séparé. Montrer que $E$ est quasi-séparé si et seulement si
  $\Top(X)$ est quasi-séparé (ou encore l'espace topologique $X$ est
  quasi-séparé, \ie l'intersection de deux ouverts quasi-compacts de
  $X$ est un ouverts quasi-compact), \emph{et} $G$ est fini ou $X$
  vide. (Comparer \hyperref[VI.exem1.33]{1.33}.)

\item Montrer que le morphisme structural $P \to e$ est quasi-compact
  si et seulement si $G$ est fini. Montrer que $E$ est cohérent (\resp
  noethérien) si et seulement si $\Top(X)$ l'est, \emph{et} $G$ est
  fini ou $X$ vide.

\item Soient $E$ un topos, $G$ un Groupe de $E$, d'où un topos
  classifiant $B_G$ (IV 2.4). Faire l'étude des conditions de finitude
  dans $B_G$, en s'inspirant de ce qui précède. Même question
  lorsqu'on part d'un pro-groupe strict $𝒢 = (G_i)_{i \in I}$ de
  $E$. En particulier, on verra que si $E$ est cohérent (\resp
  noethérien) et les $G_i$ sont cohérents, alors $B_{𝒢}$ est cohérent
  (\resp noethérien).
\end{enumerate}
\end{enonce*}

\begin{enonce*}[remark]{Exercice 2.17}\etiquette[2.17]{VI.exer2.17}
(\emph{Topos de la forme $\widehat{C}$}.)
Soient $C$ une catégorie équivalente à une catégorie $\in 𝒰$, et $E =
\widehat{C}$ la catégorie des préfaisceaux sur $C$. Par le foncteur
canonique $\varepsilon : C \to \widehat{C}$, on identifie $C$ à une
sous-catégorie pleine de $E$.
\begin{enumerate}
\renewcommand{\theenumi}{\alph {enumi}}
\renewcommand{\labelenumi}{(\theenumi)}
\item $C$ est\pageoriginale une sous-catégorie génératrice de $E$
  formée d'objets  quasi-compacts. Pour que ceux-ci soient prénoethériens (\cf \hyperref[VI.subsec1.32]{1.32}),
  il faut et il suffit que pour tout objet $X$ de $C$, tout crible de
  $X$ soit engendré par une famille finie de morphismes $X_i \to X$ de
  $C$. Pour que l'objet final de $E$ soit quasi-compact (\resp
  prénoethérien), il faut et il suffit qu'il existe une sous-catégorie
  finale de $C$ dont l'ensemble d'objets soit fini (\resp que tout
  crible de $C$ soit engendré par une sous-catégorie de $C$ dont
  l'ensemble d'objets soit fini).

\item Pour que $E$ admette une sous-catégorie génératrice formée
  d'objets cohérents (ou encore, quasi-séparés), il faut et il suffit
  que les objets de $C$ soient cohérents dans $E$, ou encore que pour
  deux flèches $Y \xrightarrow{f} X$,  $Z \xrightarrow{g} X$ dans $C$,
  l'objet $\fprod{Y}{Z}{X}$ de $E$ soit quasi-compact \ie on peut
  trouver une famille finie de carrés commutatifs dans $C$
\[
\xymatrix@=1cm{
 & T_i \ar[ld] \ar[rd]& \\
Y\ar[rd]_{f}&    & Z\ar[ld]^{g}\\
 & X  &,
}
\]
telle que  tout autre carré commutatif dans $C$
\[
\xymatrix@=1cm{
 & T \ar[ld] \ar[rd]& \\
Y\ar[rd]_{f}&    & Z\ar[ld]^{g}\\
 & X  &,
}
\]
provienne\pageoriginale d'un morphisme $T \to T_i$. (Noter que si
$(F_i)_{i \in I}$ est une famille génératrice de $E$, tout $X \in \ob C$ est isomorphe à
un facteur direct d'un des $F_i$.)

\item Pour que $E$ soit localement cohérent, \ie soit engendré par une
  famille d'objets cohérents et \emph{algébriques}, il faut et il
  suffit que les objets $X$ de $C$ soient des objets cohérents et
  algébriques de $E$, \ie que la condition de b) soit vérifiée, ainsi
  que la condition suivante: pour toute double flèche $f$, $g : Y
  \rightrightarrows Z$ de $C$ \emph{au-dessus d'un objet $X$ de} $C$,
  le noyau de cette double flèche dans $E$ est un objet quasi-compact
  de $E$, \ie il existe une famille finie de diagramme commutatifs
  dans $C$
\[
\xymatrix@C=1.5cm{
T_i \ar[r]& Y \ar@<2pt>[r]^{f,g}\ar@<-2pt>[r]&Z
}
\]
telle que tout autre diagramme commutatif dans $C$
\[
\xymatrix@C=1.5cm{
T \ar[r]& Y \ar@<2pt>[r]^{f,g}\ar@<-2pt>[r]&Z
}
\]
provienne d'un morphisme $T \to T_i$. Pour que $E$ soit algébrique,
il faut et il suffit qu'il satisfasse à la conditions de b) et à la
condition précédente, mais où on prend une double flèche
quelconque $(f, g)$ de $C$ (pas nécessairement au-dessus d'un objet
$X$ de $C$). Pour que $E$ soit quasi-séparé, il faut et il suffit
qu'il satisfasse les conditions de b), et que de plus pour deux
objets $X$, $Y\in\ob C$, le produit $X\times Y$ dans $E$ soit
quasi-compact. Pour que $E$ soit cohérent, il faut et il suffit
qu'il satisfasse aux deux conditions précédentes, et qu'il existe un
sous-catégorie finale de $C$ dont l'ensemble sous-jacent soit fini.

\item Donner\pageoriginale un exemple où les objets de la
  sous-catégorie génératrice $C$ de $E$ sont prénoethériens, mais où $E$ n'admet pas de famille
  génératrice formée d'objets cohérents (\ie (b)) où les objets de $C$
  ne sont pas tous cohérents). Prendre pour ceci pour $C$ la catégorie
  ayant des objets $X$, $T_i$ $(i = 0, 1, \ldots)$ et $e$ (l'objet
  final), les seuls morphismes entre ces objets en plus des morphismes
  structuraux dans $e$ et des identités, étant des morphismes $u_i :
  T_i \to X$, $v_i : X \to T_i$ soumis aux conditions $u_i v_i =
  \id_X$, et enfin $p_i = v_i u_i$ (satisfaisant nécessairement $p_i^2
  = \id_{T_i}$). On vérifie que les seuls cribles de $X$ ou d'un $T_i$
  sont les deux cribles triviaux, et que $e$ admet exactement un
  crible non trivial, donc en vertu de a), les objets de $C$ sont des
  objets prénoethériens de $E$. Cependant, le produit $X \times X$
  dans $E$ n'est pas quasi-compact, car il ne satisfait pas au critère
  de b).

\item Donner un exemple où la sous-catégorie génératrice $C$ de $E$
  est formée d'objets cohérents de $E$, mais où $E$ n'est pas
  localement cohérent. Prendre pour ceci pour $C$ la catégorie dont
  l'ensemble des objets est formée d'objets distincts $e$ (l'objet
  final), $Y$, $Z$ et $T_i$ $(i=0, 1, \ldots)$, avec comme seuls
  morphismes entre ces objets, en plus des morphismes dans l'objet
  final $e$ et des morphismes identiques, des morphismes $f$, $g : Y
  \rightrightarrows Y$, soumis aux conditions $fu_i = gu_i$. On
  vérifiera que $C$ satisfait à la condition de b) (avec un peu de
  patience; on trouve que tous les produits fibrés d'objets de $C$
  sont isomorphes à $∅_E$ ou sont dans $C$,\pageoriginale à
  l'exception de $\fprod{Z}{Z}{e}$ et $\fprod{Y}{Z}{e}$ qui sont recouverts par deux
  éléments de $C$), mais évidemment $\Ker(f, g)$ ne satisfait pas à la
  condition de c).

\item Donner un exemple où $C$ est stable par produits fibrés (a
  fortiori, $E$ est un topos localement cohérent \ie localement
  algébrique) mais où $E$ n'est pas un topos algébrique. (Prendre
  l'exemple donné dans d), et la sous-catégorie pleine $C'$ de $E$
  formée des objets $Y$, $Z$, $T_{i}$ $(i > 0)$ et $∅_{E}$).

\item Supposons la catégorie $C$ finie. Prouver que le topos $E$ est
  noethérien, que ses objets noethériens (\ie quasi-compacts) sont les
  contrafoncteurs $F : C^{\circ} \to (\Ens)$ tels que pour tout objet
  $X$ de $C$, $F(X)$ soit un ensemble \emph{fini}, et que les
  foncteurs fibres sur $E$ transforment objets noethériens en
  ensembles finis (\cf IV 7.6. h)).

\item Supposons que tout morphisme $f: X \to X$ de $C$ se factorise en
  un composé $ip$, où $i$ est un monomorphisme et où $p$ admet un
  inverse à droite. Soit $X$ un objet de $C$, $I(X)$ l'ensemble de ses
  sous-objets au sens de $C$, considéré comme un ensemble ordonné,
  $S(I(X))$ l'ensemble des parties $U$ de $I(X)$ telles que pour deux
  éléments $X'$, $X'' \in I(X)$, $X'\in U$ et $X''\leq X'$ implique
  $X''\in U$. Montrer que l'ensemble ordonné des sous-objets dans $E$
  de $X$ est isomorphe à l'ensemble $S(I(X))$ (ordonné par
  inclusion). En conclure que si $I(X)$ est fini, alors $X$ est un
  objet prénoethérien de $E$. En particulier, si (en plus de la
  conditions de factorisation ci-dessus) $C$ est stable par produits
  fibrés (\resp par limites projectives finies) et\pageoriginale si
  pour tout $X \in
  \ob C$, l'ensemble $I(X)$ des sous-objets de $X$ dans $C$ est fini,
  alors le topos $E$ est localement noethérien (\resp noethérien).

\item Soit $C$ la catégorie des ensembles finis (ou, au choix, des
  ensembles finis non vides) $\in 𝒰$, de sorte que $E = \widehat{C}$
  est la catégorie des ensembles simpliciaux augmentés (\resp des
  ensembles simpliciaux tout court). Montrer que $E$ est un topos
  noethérien. (Utiliser h).) Prenant un pro-objet $(X_i)_{i \in I}$
  non essentiellement constant de $C$, montrer que $E$ admet des
  foncteurs fibres qui ne transforment pas objets noethériens en
  objets noethériens (\ie en ensembles finis).

\item On considère le diagramme d'implications suivant de propriétés
  pour un topos $E$:
\[
\xymatrix@R=.75cm{
**[r]\txt{E est noethérien} \ar@{=>}[d] \ar@{=>}[r]& **[r]\txt{E est
  localement  noethérien} \ar@{=>}[d] & \ar@{=>}[rdd]\\
**[r]\txt{E est cohérent} \ar@{=>}[r] & **[r]\txt{E
est algébrique} \ar@{=>}[d]\\
 & **[r]\txt<9.5pc>{E est loc. cohérent, i.e.\\E \raisebox{-1ex}{cohalg}
  engendre E} \ar@{=>}[d]& & \txt{E \raisebox{-1ex}{prénoeth} engendre
  E} \ar@{=>}[ldd] \\
 & **[r]\txt{E \raisebox{-1ex}{coh} engendre E} \ar@{=>}[d]&\\
 & **[r]\txt{E \raisebox{-1ex}{qucpct} engendre E}  &&.\\
}
\]
Montrer que toutes les implications de ce diagramme sont strictes,
et qu'il n'y a pas entre les notions envisagées d'autres implications
que les implications composées du diagramme précédent. Ici $E_{\coh},
E_{\text{qucpct}}$, $E_{\text{prénoeth}}$,
$E_{\text{cohalg}}$\pageoriginale
désignent respectivement les sous-catégories pleines de $E$ formées des objets cohérents, \resp
quasi-compacts, \resp prénoethériens, \resp cohérents et algébriques de
$E$. (On se bornera à des topos de la forme $\widehat{C}$, en
utilisant les résultats énoncés dans d), e), f).)

\item Résoudre la question suivante (dont le rédacteur de ces lignes
  ignore la réponse) : $E = \widehat{C}$ peut-il être localement
  noethérien sans que les $\Hom(X, Y)$ (pour $X$, $Y \in \ob C$) soient
  finis ?
\end{enumerate}
\end{enonce*}

\begin{enonce*}[remark]{Exercice 2.18}\etiquette[2.18]{VI.exer2.18}
Soit $E$ un topos.
\begin{enumerate}
\renewcommand{\theenumi}{\alph{enumi}}
\renewcommand{\labelenumi}{\theenumi)}
\item Soit $X$ un objet prénoethérien de $E$. Montrer que $X$ est
  isomorphe à une somme finie d'objets connexes (IV 4.3.5) de
  $E$. (Montrer par l'absurde qu'une suite croissante de partitions de
  $X$ (IV 8.7) est stationnaire.)

\item Soient $X$ un objet de $E$, $(f_i : X_i \to X)_{i \in I}$ une
  famille couvrante de $X$. Montrer que si chacun des $X_i$ est
  isomorphe à une somme d'objets connexes de $E$, il en est de même de
  $X$. (Se ramener au cas où les $X_i$ sont connexes, puis au cas où
  ce sont des sous-objets de $X$, et considérer alors sur l'ensemble
  d'indices $I$ la relation d'équivalence $R$ engendrée par la
  relation $X_i \cap X_j \neq \phi_E$, et montrer que si $J = I/R$,
  $X$ est somme des $X(i) = \displaystyle{\mathop{\Sup}_{i \in I}} X_i
  .)$

\item Conclure de b) que si on désigne par $(X_i)_{i \in I}$ une
  famille génératrice de $E$, alors $E$ est localement connexe (IV
  8.7. $\ell$)) si\pageoriginale et seulement si chacun des $X_i$ $(i \in I)$
  est isomorphe à une somme d'objets connexes de $E$.

\item Conclure de a) et c) que si $E$ admet une sous-catégorie
  génératrice formée d'objets prénoethériens, en particulier si $E$
  est localement noethérien, alors $E$ est localement connexe, et a
  fortiori est isomorphe au topos somme (IV 8.7 b)) d'une famille de
  topos connexes (IV 8.7 e)) \ie dont l'objet final est connexe. (En
  particulier, on peut associer à $E$, pour tout morphisme $f : P \to
  E$, où $P$ est un topos « connexe non vide et simplement
  connexe » (IV 2.7.5) un pro-groupe fondamental $\pi_1(E, f)$; on
  notera que si $E$ est localement noethérien « non vide » on peut
  toujours trouver un tel $f$, avec $P$ le topos ponctuel, grâce à
  \textsc{Deligne} ([9]).
\end{enumerate}
\end{enonce*}

